% F02：本书原创矢量示意；依据所在节的定义、证明或明确模型。
\begin{figure}[htbp]
\centering
\begin{tikzpicture}[x=1cm,y=1cm,>=stealth,every node/.style={font=\small},line width=.65pt]
\draw[->,gray] (-0.12,0)--(3.4,0) node[right] {\(x\)};\draw[->,gray] (0.0,-0.12)--(0.0,2.8);
\node[] at (1.65,3.65) {集中};
\draw[AnalysisBlue,line width=.8pt] (0.0,0)--(0.0,0.6)--(1.5,0.6)--(1.5,0);
\draw[orange!80!black,line width=.8pt] (0.0,0)--(0.0,1.2)--(0.75,1.2)--(0.75,0);
\draw[black,line width=.8pt] (0.0,0)--(0.0,2.4)--(0.375,2.4)--(0.375,0);
\node[] at (1.65,-0.55) {\(n\cdot\mathbf1_{(0,1/n)}\)};
\node[font=\footnotesize] at (1.65,3.05) {\(n=2,4,8\)};
\draw[->,gray] (4.28,0)--(7.800000000000001,0) node[right] {\(x\)};\draw[->,gray] (4.4,-0.12)--(4.4,2.8);
\node[] at (6.050000000000001,3.65) {逃逸};
\draw[AnalysisBlue,line width=.8pt] (5.0600000000000005,0)--(5.0600000000000005,1.5)--(5.390000000000001,1.5)--(5.390000000000001,0);
\draw[orange!80!black,line width=.8pt] (5.720000000000001,0)--(5.720000000000001,1.5)--(6.050000000000001,1.5)--(6.050000000000001,0);
\draw[black,line width=.8pt] (7.040000000000001,0)--(7.040000000000001,1.5)--(7.370000000000001,1.5)--(7.370000000000001,0);
\node[] at (6.050000000000001,-0.55) {\(\mathbf1_{[n,n+1]}\)};
\node[font=\footnotesize] at (6.050000000000001,3.05) {\(n=2,4,8\)};
\draw[->,gray] (8.680000000000001,0)--(12.200000000000001,0) node[right] {\(x\)};\draw[->,gray] (8.8,-0.12)--(8.8,2.8);
\node[] at (10.450000000000001,3.65) {总量增大};
\draw[AnalysisBlue,line width=.8pt] (8.8,0)--(8.8,1.5)--(9.540000000000001,1.5)--(9.540000000000001,0);
\draw[orange!80!black,line width=.8pt] (8.8,0)--(8.8,1.5)--(10.280000000000001,1.5)--(10.280000000000001,0);
\draw[black,line width=.8pt] (8.8,0)--(8.8,1.5)--(11.760000000000002,1.5)--(11.760000000000002,0);
\node[] at (10.450000000000001,-0.55) {\(\mathbf1_{[0,n]}\)};
\node[font=\footnotesize] at (10.450000000000001,3.05) {\(n=2,4,8\)};

\end{tikzpicture}
\caption[集中、逃逸与幅值尾部]{集中、逃逸与幅值尾部。集中族的积分恒为一且不一致可积；后两族按本书幅值尾部定义均一致可积，但分别存在空间逃逸和积分总量无界。各面板坐标尺度不同。}
\label{b1:fig:visual-f02}
\end{figure}
